\documentclass[12pt, a4paper]{article}

% ── Encoding & Language ──────────────────────────────────────────────────────
\usepackage[utf8]{inputenc}
\usepackage[T1]{fontenc}
\usepackage[spanish]{babel}

% ── Geometry & Layout ────────────────────────────────────────────────────────
\usepackage[top=2.5cm, bottom=2.5cm, left=2cm, right=2cm]{geometry}
\usepackage{pdflscape}

% ── Typography ───────────────────────────────────────────────────────────────
\usepackage{lmodern}
\usepackage{microtype}
\usepackage{setspace}
\onehalfspacing

% ── Color ────────────────────────────────────────────────────────────────────
\usepackage[table,xcdraw]{xcolor}
\definecolor{espe-blue}{RGB}{0, 51, 102}
\definecolor{espe-gold}{RGB}{204, 153, 0}
\definecolor{header-bg}{RGB}{0, 51, 102}
\definecolor{row-alt}{RGB}{230, 240, 255}
\definecolor{critica}{RGB}{220, 53, 69}
\definecolor{alta}{RGB}{255, 140, 0}
\definecolor{media}{RGB}{40, 167, 69}
\definecolor{paso-green}{RGB}{40, 167, 69}
\definecolor{pendiente-orange}{RGB}{255, 140, 0}

% ── Tables ───────────────────────────────────────────────────────────────────
\usepackage{booktabs}
\usepackage{array}
\usepackage{longtable}
\usepackage{amssymb}

% ── Graphics & Boxes ─────────────────────────────────────────────────────────
\usepackage{graphicx}
\usepackage[most,breakable]{tcolorbox}
\usepackage{tikz}

% ── Lists ────────────────────────────────────────────────────────────────────
\usepackage{enumitem}

% ── Header & Footer ──────────────────────────────────────────────────────────
\usepackage{fancyhdr}
\pagestyle{fancy}
\fancyhf{}
\fancyhead[L]{\small\textcolor{espe-blue}{\textbf{SpeechNotes -- SQA}}}
\fancyhead[R]{\small\textcolor{espe-blue}{Persona D -- QA Manual}}
\fancyfoot[C]{\small\thepage}
\fancyfoot[R]{\small\textcolor{gray}{Sprint 18--21 Jul 2026}}
\renewcommand{\headrulewidth}{1.5pt}
\renewcommand{\headrule}{\hbox to\headwidth{\color{espe-blue}\leaders\hrule height \headrulewidth\hfill}}

% ── Hyperlinks ───────────────────────────────────────────────────────────────
\usepackage[colorlinks=true, linkcolor=espe-blue, urlcolor=espe-blue,
            pdftitle={SQA - Matriz de Rastreabilidad - SpeechNotes},
            pdfauthor={Persona D - QA Manual}]{hyperref}

% ── Helpers ──────────────────────────────────────────────────────────────────
\newcommand{\prioChip}[2]{%
  \colorbox{#1}{\textcolor{white}{\footnotesize\textbf{#2}}}%
}
\newcommand{\statusChip}[2]{%
  \colorbox{#1}{\textcolor{white}{\footnotesize\textbf{#2}}}%
}
\newcolumntype{L}[1]{>{\raggedright\arraybackslash}p{#1}}
\newcolumntype{C}[1]{>{\centering\arraybackslash}p{#1}}

% =============================================================================
\begin{document}
\thispagestyle{empty}

% ── PORTADA ──────────────────────────────────────────────────────────────────
\begin{tikzpicture}[remember picture, overlay]
  \fill[espe-blue] (current page.north west) rectangle ([yshift=-6cm]current page.north east);
  \fill[espe-gold]  ([yshift=-6cm]current page.north west) rectangle ([yshift=-6.4cm]current page.north east);
\end{tikzpicture}

\vspace*{1.8cm}
\begin{center}
  {\color{white}\fontsize{22}{26}\selectfont\textbf{PLAN DE ASEGURAMIENTO DE CALIDAD}}\\[0.3cm]
  {\color{white}\fontsize{16}{20}\selectfont\textit{SpeechNotes -- Sprint 18 al 21 de Julio 2026}}\\[0.5cm]
  {\color{espe-gold}\fontsize{18}{22}\selectfont\textbf{Persona D: QA Manual \& Dise\~no de Pruebas}}\\[1.5cm]
  {\color{white}\large Matriz de Rastreabilidad de Requisitos}\\
  {\color{white}\large y Casos de Prueba Manuales}\\[2cm]
\end{center}

\begin{tcolorbox}[
  enhanced,
  colback=white, colframe=espe-gold, boxrule=1.5pt, arc=6pt,
  width=12cm, center, fontupper=\small
]
  \begin{tabular}{ll}
    \textbf{Proyecto:}    & SpeechNotes \\
    \textbf{Documento:}   & Secci\'on B.a -- B.b \\
    \textbf{Responsable:} & Persona D -- QA Manual \\
    \textbf{Fecha:}       & 18 de julio de 2026 \\
    \textbf{Rama:}        & \texttt{planning-SQA} \\
    \textbf{Frontend:}    & \texttt{http://localhost:3006} \\
    \textbf{Backend:}     & \texttt{http://localhost:9443} \\
    \textbf{Entrega:}     & 20 de julio de 2026 \\
  \end{tabular}
\end{tcolorbox}

\newpage

% ── TOC ──────────────────────────────────────────────────────────────────────
\tableofcontents
\newpage

% ── 1. ALCANCE ───────────────────────────────────────────────────────────────
\section{Alcance del QA Manual}

\begin{tcolorbox}[enhanced, breakable, colback=row-alt, colframe=espe-blue,
  boxrule=1pt, arc=4pt,
  title={\textbf{Flujos incluidos en la cobertura manual}}, fonttitle=\color{white},
  attach boxed title to top left={yshift=-2mm, xshift=4mm},
  boxed title style={colback=espe-blue}]
\begin{itemize}[noitemsep]
  \item Login con usuario demo (\texttt{demo@speechnotes.app})
  \item Prueba de micr\'ofono
  \item Grabaci\'on en vivo e inicio/detenci\'on
  \item Transcripci\'on en vivo (Realtime ASR)
  \item Upload de archivo de audio (.mp3 / .wav)
  \item Consulta de transcripciones guardadas
  \item Edici\'on y guardado de transcripciones
  \item B\'usqueda dentro de transcripciones
  \item Formateo con Inteligencia Artificial
  \item Chat contextual sobre una nota
\end{itemize}
\end{tcolorbox}

\vspace{0.5cm}

\begin{tcolorbox}[enhanced, colback=white, colframe=critica,
  boxrule=1pt, arc=4pt,
  title={\textbf{Fuera de alcance}}, fonttitle=\color{white},
  attach boxed title to top left={yshift=-2mm, xshift=4mm},
  boxed title style={colback=critica}]
\begin{itemize}[noitemsep]
  \item Implementar automatizaci\'on (responsabilidad de Personas B y C)
  \item Corregir c\'odigo fuente
  \item Configurar SonarCloud o CI/CD
  \item Pruebas formales de carga (riesgo documentado)
\end{itemize}
\end{tcolorbox}

\newpage

% ── 2. MATRIZ DE RASTREABILIDAD ──────────────────────────────────────────────
\section{B.a Matriz de Rastreabilidad de Requisitos}

La siguiente tabla traza cada Requisito Funcional (RF) al Caso de Prueba (CP) que lo
valida, indicando prioridad, m\'odulo afectado y criterio de aceptaci\'on definido.

\vspace{0.4cm}

\begin{landscape}
\begin{longtable}{C{1.4cm} L{4.8cm} C{1.5cm} C{2.2cm} C{2.8cm} L{6.5cm}}
\toprule
\rowcolor{header-bg}
\textcolor{white}{\textbf{ID Req}} &
\textcolor{white}{\textbf{Requisito Funcional}} &
\textcolor{white}{\textbf{Caso}} &
\textcolor{white}{\textbf{Prioridad}} &
\textcolor{white}{\textbf{M\'odulo}} &
\textcolor{white}{\textbf{Criterio de Aceptaci\'on}} \\
\midrule
\endfirsthead
\toprule
\rowcolor{header-bg}
\textcolor{white}{\textbf{ID Req}} &
\textcolor{white}{\textbf{Requisito Funcional}} &
\textcolor{white}{\textbf{Caso}} &
\textcolor{white}{\textbf{Prioridad}} &
\textcolor{white}{\textbf{M\'odulo}} &
\textcolor{white}{\textbf{Criterio de Aceptaci\'on}} \\
\midrule
\endhead
\midrule\multicolumn{6}{r}{\small\textit{Contin\'ua\ldots}}\\\endfoot
\bottomrule\endlastfoot

RF-01 & El usuario puede iniciar sesi\'on & CP-01 & \prioChip{alta}{Alta} & Auth &
  El usuario demo accede al dashboard sin error de login. \\[4pt]

\rowcolor{row-alt}
RF-02 & El usuario puede probar el micr\'ofono & CP-02 & \prioChip{alta}{Alta} & Audio / Frontend &
  El panel de prueba muestra actividad cuando el usuario habla. \\[4pt]

RF-03 & El usuario puede iniciar/detener grabaci\'on & CP-03 & \prioChip{critica}{Cr\'itica} & Audio / Backend &
  La UI cambia a \textit{grabando}, el contador avanza y al detener se procesa la grabaci\'on. \\[4pt]

\rowcolor{row-alt}
RF-04 & El sistema muestra transcripci\'on en vivo & CP-04 & \prioChip{critica}{Cr\'itica} & Realtime ASR &
  El panel en vivo agrega segmentos durante la grabaci\'on sin bucles incoherentes. \\[4pt]

RF-05 & El usuario puede subir un audio para transcribir & CP-05 & \prioChip{alta}{Alta} & Upload / Backend &
  El archivo se procesa y aparece una transcripci\'on nueva en la lista. \\[4pt]

\rowcolor{row-alt}
RF-06 & El usuario puede consultar transcripciones guardadas & CP-06 & \prioChip{alta}{Alta} & Transcripciones &
  La lista carga documentos y el visor muestra el contenido seleccionado sin timeout. \\[4pt]

RF-07 & El usuario puede editar y guardar una transcripci\'on & CP-07 & \prioChip{media}{Media} & Editor &
  Los cambios persisten despu\'es de recargar; el Markdown no se corrompe. \\[4pt]

\rowcolor{row-alt}
RF-08 & El usuario puede buscar dentro de transcripciones & CP-08 & \prioChip{media}{Media} & B\'usqueda &
  Una palabra existente retorna resultados; una inexistente muestra estado vac\'io sin error. \\[4pt]

RF-09 & El usuario puede formatear una transcripci\'on con IA & CP-09 & \prioChip{media}{Media} & IA / Formatter &
  El job inicia, la UI muestra estado y al completar el contenido queda formateado. \\[4pt]

\rowcolor{row-alt}
RF-10 & El usuario puede usar chat contextual sobre una nota & CP-10 & \prioChip{media}{Media} & Chat / RAG &
  El chat responde sin error HTTP usando informaci\'on de la transcripci\'on; el historial es visible. \\[4pt]

\end{longtable}

\vspace{0.3cm}
\begin{tcolorbox}[colback=row-alt, colframe=espe-blue, boxrule=0.8pt, arc=3pt]
\footnotesize
\textbf{Leyenda:}\quad
\prioChip{critica}{Cr\'itica} = flujo bloqueante\quad
\prioChip{alta}{Alta} = funcionalidad central\quad
\prioChip{media}{Media} = valor a\~nadido
\end{tcolorbox}
\end{landscape}

\newpage

% ── 3. CASOS DE PRUEBA ───────────────────────────────────────────────────────
\section{B.b Casos de Prueba Manuales}

Cada caso incluye: precondiciones, pasos de ejecuci\'on, resultado esperado, resultado
obtenido y estado final.

% ── CP-01 ────────────────────────────────────────────────────────────────────
\begin{tcolorbox}[enhanced, breakable, colback=white, colframe=espe-blue,
  boxrule=1pt, arc=5pt,
  title={\textbf{CP-01} \enspace Login con usuario demo \hfill \prioChip{alta}{Alta} \enspace \statusChip{paso-green}{PASO}},
  fonttitle=\color{white},
  attach boxed title to top left={yshift=-2mm,xshift=4mm},
  boxed title style={colback=espe-blue,arc=3pt}]
\begin{minipage}[t]{0.48\linewidth}
\textbf{Precondiciones:}
\begin{itemize}[noitemsep,topsep=2pt]
  \item Frontend corriendo en puerto 3006.
  \item Backend corriendo en puerto 9443.
  \item Base de datos disponible.
\end{itemize}
\textbf{Pasos de Ejecuci\'on:}
\begin{enumerate}[noitemsep,topsep=2pt]
  \item Abrir \texttt{http://localhost:3006/login}.
  \item Ingresar \texttt{demo@speechnotes.app}.
  \item Ingresar \texttt{demo123}.
  \item Presionar el bot\'on de inicio de sesi\'on.
  \item Verificar redirecci\'on al dashboard.
\end{enumerate}
\end{minipage}\hfill
\begin{minipage}[t]{0.48\linewidth}
\textbf{Resultado Esperado:}
\begin{itemize}[noitemsep,topsep=2pt]
  \item El usuario llega a \texttt{/dashboard}.
  \item No aparece \texttt{OAuthSignin}.
  \item No aparece error de credenciales.
\end{itemize}
\vspace{4pt}
\textbf{Resultado Obtenido:}
\begin{itemize}[noitemsep,topsep=2pt]
  \item[$\checkmark$] El usuario llega a \texttt{/dashboard}.
  \item[$\checkmark$] Sin \texttt{OAuthSignin}.
  \item[$\checkmark$] Sin error de credenciales.
\end{itemize}
\end{minipage}
\end{tcolorbox}

\vspace{0.3cm}

% ── CP-02 ────────────────────────────────────────────────────────────────────
\begin{tcolorbox}[enhanced, breakable, colback=white, colframe=espe-blue,
  boxrule=1pt, arc=5pt,
  title={\textbf{CP-02} \enspace Prueba de micr\'ofono \hfill \prioChip{alta}{Alta} \enspace \statusChip{paso-green}{PASO}},
  fonttitle=\color{white},
  attach boxed title to top left={yshift=-2mm,xshift=4mm},
  boxed title style={colback=espe-blue,arc=3pt}]
\begin{minipage}[t]{0.48\linewidth}
\textbf{Precondiciones:}
\begin{itemize}[noitemsep,topsep=2pt]
  \item Usuario autenticado.
  \item Micr\'ofono conectado.
  \item Permisos del navegador disponibles.
\end{itemize}
\textbf{Pasos de Ejecuci\'on:}
\begin{enumerate}[noitemsep,topsep=2pt]
  \item Entrar a \texttt{/dashboard}.
  \item Abrir la herramienta de prueba de micr\'ofono.
  \item Aceptar permisos del navegador.
  \item Hablar durante 5 a 10 segundos.
  \item Verificar indicador de nivel/pico.
  \item Detener la prueba.
\end{enumerate}
\end{minipage}\hfill
\begin{minipage}[t]{0.48\linewidth}
\textbf{Resultado Esperado:}
\begin{itemize}[noitemsep,topsep=2pt]
  \item El indicador visual cambia al hablar.
  \item La prueba se detiene sin bloquear el micr\'ofono.
\end{itemize}
\vspace{4pt}
\textbf{Resultado Obtenido:}
\begin{itemize}[noitemsep,topsep=2pt]
  \item[$\checkmark$] El indicador visual cambia al hablar.
  \item[$\checkmark$] Sin bloqueo del micr\'ofono.
\end{itemize}
\end{minipage}
\end{tcolorbox}

\vspace{0.3cm}

% ── CP-03 ────────────────────────────────────────────────────────────────────
\begin{tcolorbox}[enhanced, breakable, colback=white, colframe=critica,
  boxrule=1.5pt, arc=5pt,
  title={\textbf{CP-03} \enspace Iniciar y detener grabaci\'on \hfill \prioChip{critica}{Cr\'itica} \enspace \statusChip{paso-green}{PASO}},
  fonttitle=\color{white},
  attach boxed title to top left={yshift=-2mm,xshift=4mm},
  boxed title style={colback=critica,arc=3pt}]
\begin{minipage}[t]{0.48\linewidth}
\textbf{Precondiciones:}
\begin{itemize}[noitemsep,topsep=2pt]
  \item Backend corriendo en puerto 9443.
  \item Frontend corriendo en puerto 3006.
  \item Usuario autenticado.
  \item Micr\'ofono disponible.
\end{itemize}
\textbf{Pasos de Ejecuci\'on:}
\begin{enumerate}[noitemsep,topsep=2pt]
  \item Entrar a \texttt{/dashboard}.
  \item Presionar el bot\'on de micr\'ofono.
  \item Aceptar permisos del navegador.
  \item Hablar durante 10 segundos.
  \item Presionar detener.
  \item Confirmar la detenci\'on.
\end{enumerate}
\end{minipage}\hfill
\begin{minipage}[t]{0.48\linewidth}
\textbf{Resultado Esperado:}
\begin{itemize}[noitemsep,topsep=2pt]
  \item El estado cambia a \textit{grabando}.
  \item El contador aumenta.
  \item Al detener, aparece estado de procesamiento.
  \item La grabaci\'on se agrega a la lista de transcripciones.
\end{itemize}
\vspace{4pt}
\textbf{Resultado Obtenido:}
\begin{itemize}[noitemsep,topsep=2pt]
  \item[$\checkmark$] Estado \textit{grabando} visible.
  \item[$\checkmark$] Contador aumenta.
  \item[$\checkmark$] Estado de procesamiento al detener.
  \item[$\checkmark$] Transcripci\'on agregada a la lista.
\end{itemize}
\end{minipage}
\end{tcolorbox}

\newpage

% ── CP-04 ────────────────────────────────────────────────────────────────────
\begin{tcolorbox}[enhanced, breakable, colback=white, colframe=critica,
  boxrule=1.5pt, arc=5pt,
  title={\textbf{CP-04} \enspace Transcripci\'on en vivo \hfill \prioChip{critica}{Cr\'itica} \enspace \statusChip{paso-green}{PASO}},
  fonttitle=\color{white},
  attach boxed title to top left={yshift=-2mm,xshift=4mm},
  boxed title style={colback=critica,arc=3pt}]
\begin{minipage}[t]{0.48\linewidth}
\textbf{Precondiciones:}
\begin{itemize}[noitemsep,topsep=2pt]
  \item Usuario autenticado.
  \item Backend ASR configurado.
  \item Socket.IO conectado.
  \item Micr\'ofono con audio claro.
\end{itemize}
\textbf{Pasos de Ejecuci\'on:}
\begin{enumerate}[noitemsep,topsep=2pt]
  \item Entrar a \texttt{/dashboard}.
  \item Iniciar grabaci\'on.
  \item Hablar frases claras durante 30--60 segundos.
  \item Observar el panel de transcripci\'on en vivo.
  \item Confirmar el aumento del contador de segmentos.
  \item Detener la grabaci\'on.
\end{enumerate}
\end{minipage}\hfill
\begin{minipage}[t]{0.48\linewidth}
\textbf{Resultado Esperado:}
\begin{itemize}[noitemsep,topsep=2pt]
  \item El panel muestra segmentos sin esperar al final.
  \item El contador de segmentos aumenta.
  \item No aparecen repeticiones incoherentes en bucle.
\end{itemize}
\vspace{4pt}
\textbf{Resultado Obtenido:}
\begin{itemize}[noitemsep,topsep=2pt]
  \item[$\checkmark$] Segmentos visibles progresivamente.
  \item[$\checkmark$] Contador de segmentos aumenta.
  \item[$\checkmark$] Sin repeticiones en bucle.
\end{itemize}
\end{minipage}
\end{tcolorbox}

\vspace{0.3cm}

% ── CP-05 ────────────────────────────────────────────────────────────────────
\begin{tcolorbox}[enhanced, breakable, colback=white, colframe=espe-blue,
  boxrule=1pt, arc=5pt,
  title={\textbf{CP-05} \enspace Upload de audio para transcribir \hfill \prioChip{alta}{Alta} \enspace \statusChip{paso-green}{PASO}},
  fonttitle=\color{white},
  attach boxed title to top left={yshift=-2mm,xshift=4mm},
  boxed title style={colback=espe-blue,arc=3pt}]
\begin{minipage}[t]{0.48\linewidth}
\textbf{Precondiciones:}
\begin{itemize}[noitemsep,topsep=2pt]
  \item Usuario autenticado.
  \item Backend disponible.
  \item Archivo \texttt{.mp3} o \texttt{.wav} de prueba disponible.
\end{itemize}
\textbf{Pasos de Ejecuci\'on:}
\begin{enumerate}[noitemsep,topsep=2pt]
  \item Entrar a \texttt{/dashboard}.
  \item Abrir la herramienta Upload.
  \item Seleccionar un archivo \texttt{.mp3} o \texttt{.wav}.
  \item Iniciar la transcripci\'on.
  \item Esperar el procesamiento.
  \item Revisar la lista de transcripciones.
\end{enumerate}
\end{minipage}\hfill
\begin{minipage}[t]{0.48\linewidth}
\textbf{Resultado Esperado:}
\begin{itemize}[noitemsep,topsep=2pt]
  \item El archivo se acepta.
  \item El sistema muestra estado de procesamiento.
  \item Al finalizar, la transcripci\'on aparece en la lista.
  \item El visor abre el texto generado.
\end{itemize}
\vspace{4pt}
\textbf{Resultado Obtenido:}
\begin{itemize}[noitemsep,topsep=2pt]
  \item[$\checkmark$] Archivo aceptado.
  \item[$\checkmark$] Estado de procesamiento visible.
  \item[$\checkmark$] Transcripci\'on en lista.
  \item[$\checkmark$] Visor abre el texto.
\end{itemize}
\end{minipage}
\end{tcolorbox}

\vspace{0.3cm}

% ── CP-06 ────────────────────────────────────────────────────────────────────
\begin{tcolorbox}[enhanced, breakable, colback=white, colframe=espe-blue,
  boxrule=1pt, arc=5pt,
  title={\textbf{CP-06} \enspace Consultar transcripciones guardadas \hfill \prioChip{alta}{Alta} \enspace \statusChip{paso-green}{PASO}},
  fonttitle=\color{white},
  attach boxed title to top left={yshift=-2mm,xshift=4mm},
  boxed title style={colback=espe-blue,arc=3pt}]
\begin{minipage}[t]{0.48\linewidth}
\textbf{Precondiciones:}
\begin{itemize}[noitemsep,topsep=2pt]
  \item Usuario autenticado.
  \item Al menos una transcripci\'on guardada.
\end{itemize}
\textbf{Pasos de Ejecuci\'on:}
\begin{enumerate}[noitemsep,topsep=2pt]
  \item Entrar a \texttt{/dashboard}.
  \item Verificar que la lista de transcripciones cargue.
  \item Seleccionar la transcripci\'on m\'as reciente.
  \item Navegar a una transcripci\'on anterior.
  \item Volver a la transcripci\'on m\'as reciente.
\end{enumerate}
\end{minipage}\hfill
\begin{minipage}[t]{0.48\linewidth}
\textbf{Resultado Esperado:}
\begin{itemize}[noitemsep,topsep=2pt]
  \item La lista muestra transcripciones existentes.
  \item El visor cambia seg\'un la nota seleccionada.
  \item Sin timeout al cargar documentos.
\end{itemize}
\vspace{4pt}
\textbf{Resultado Obtenido:}
\begin{itemize}[noitemsep,topsep=2pt]
  \item[$\checkmark$] Lista con transcripciones.
  \item[$\checkmark$] Visor actualiza contenido.
  \item[$\checkmark$] Sin timeout.
\end{itemize}
\end{minipage}
\end{tcolorbox}

\newpage

% ── CP-07 ────────────────────────────────────────────────────────────────────
\begin{tcolorbox}[enhanced, breakable, colback=white, colframe=espe-blue,
  boxrule=1pt, arc=5pt,
  title={\textbf{CP-07} \enspace Editar y guardar una transcripci\'on \hfill \prioChip{media}{Media} \enspace \statusChip{paso-green}{PASO}},
  fonttitle=\color{white},
  attach boxed title to top left={yshift=-2mm,xshift=4mm},
  boxed title style={colback=espe-blue,arc=3pt}]
\begin{minipage}[t]{0.48\linewidth}
\textbf{Precondiciones:}
\begin{itemize}[noitemsep,topsep=2pt]
  \item Usuario autenticado.
  \item Debe existir una transcripci\'on guardada.
\end{itemize}
\textbf{Pasos de Ejecuci\'on:}
\begin{enumerate}[noitemsep,topsep=2pt]
  \item Abrir una transcripci\'on en el visor.
  \item Activar edici\'on si aplica.
  \item Agregar una l\'inea de prueba al final.
  \item Guardar cambios.
  \item Recargar la p\'agina.
  \item Abrir la misma transcripci\'on.
\end{enumerate}
\end{minipage}\hfill
\begin{minipage}[t]{0.48\linewidth}
\textbf{Resultado Esperado:}
\begin{itemize}[noitemsep,topsep=2pt]
  \item El cambio se guarda sin error.
  \item La l\'inea persiste despu\'es de recargar.
  \item El Markdown no se corrompe.
\end{itemize}
\vspace{4pt}
\textbf{Resultado Obtenido:}
\begin{itemize}[noitemsep,topsep=2pt]
  \item[$\checkmark$] Cambio guardado sin error.
  \item[$\checkmark$] L\'inea persiste tras recarga.
  \item[$\checkmark$] Markdown sin corrupci\'on.
\end{itemize}
\end{minipage}
\end{tcolorbox}

\vspace{0.3cm}

% ── CP-08 ────────────────────────────────────────────────────────────────────
\begin{tcolorbox}[enhanced, breakable, colback=white, colframe=espe-blue,
  boxrule=1pt, arc=5pt,
  title={\textbf{CP-08} \enspace Buscar dentro de transcripciones \hfill \prioChip{media}{Media} \enspace \statusChip{paso-green}{PASO}},
  fonttitle=\color{white},
  attach boxed title to top left={yshift=-2mm,xshift=4mm},
  boxed title style={colback=espe-blue,arc=3pt}]
\begin{minipage}[t]{0.48\linewidth}
\textbf{Precondiciones:}
\begin{itemize}[noitemsep,topsep=2pt]
  \item Usuario autenticado.
  \item Transcripci\'on con contenido conocido.
\end{itemize}
\textbf{Pasos de Ejecuci\'on:}
\begin{enumerate}[noitemsep,topsep=2pt]
  \item Entrar a \texttt{/dashboard}.
  \item Abrir la b\'usqueda global.
  \item Buscar una palabra existente.
  \item Abrir el resultado.
  \item Buscar una palabra inexistente.
\end{enumerate}
\end{minipage}\hfill
\begin{minipage}[t]{0.48\linewidth}
\textbf{Resultado Esperado:}
\begin{itemize}[noitemsep,topsep=2pt]
  \item La palabra existente devuelve resultados.
  \item El resultado abre la nota correcta.
  \item La inexistente muestra estado vac\'io sin error.
\end{itemize}
\vspace{4pt}
\textbf{Resultado Obtenido:}
\begin{itemize}[noitemsep,topsep=2pt]
  \item[$\checkmark$] Resultados correctos para palabra existente.
  \item[$\checkmark$] Nota correcta abierta.
  \item[$\checkmark$] Estado vac\'io para inexistente.
\end{itemize}
\end{minipage}
\end{tcolorbox}

\vspace{0.3cm}

% ── CP-09 ────────────────────────────────────────────────────────────────────
\begin{tcolorbox}[enhanced, breakable, colback=white, colframe=espe-blue,
  boxrule=1pt, arc=5pt,
  title={\textbf{CP-09} \enspace Formateo IA de una transcripci\'on \hfill \prioChip{media}{Media} \enspace \statusChip{pendiente-orange}{PENDIENTE}},
  fonttitle=\color{white},
  attach boxed title to top left={yshift=-2mm,xshift=4mm},
  boxed title style={colback=espe-blue,arc=3pt}]
\begin{minipage}[t]{0.48\linewidth}
\textbf{Precondiciones:}
\begin{itemize}[noitemsep,topsep=2pt]
  \item Usuario autenticado.
  \item Transcripci\'on seleccionada.
  \item Credenciales IA configuradas localmente.
\end{itemize}
\textbf{Pasos de Ejecuci\'on:}
\begin{enumerate}[noitemsep,topsep=2pt]
  \item Abrir una transcripci\'on.
  \item Presionar la acci\'on de formateo IA.
  \item Iniciar el proceso.
  \item Observar el estado de procesamiento.
  \item Esperar finalizaci\'on.
  \item Revisar el contenido actualizado.
\end{enumerate}
\end{minipage}\hfill
\begin{minipage}[t]{0.48\linewidth}
\textbf{Resultado Esperado:}
\begin{itemize}[noitemsep,topsep=2pt]
  \item El job de formateo inicia.
  \item La UI muestra estado mientras procesa.
  \item Al completar, el contenido aparece formateado.
  \item No se pierde el contenido original.
\end{itemize}
\vspace{4pt}
\textbf{Resultado Obtenido:}
\begin{itemize}[noitemsep,topsep=2pt]
  \item \textit{Pendiente de ejecuci\'on.}
  \item Requiere credenciales IA locales.
\end{itemize}
\end{minipage}
\end{tcolorbox}

\newpage

% ── CP-10 ────────────────────────────────────────────────────────────────────
\begin{tcolorbox}[enhanced, breakable, colback=white, colframe=espe-blue,
  boxrule=1pt, arc=5pt,
  title={\textbf{CP-10} \enspace Chat contextual sobre una nota \hfill \prioChip{media}{Media} \enspace \statusChip{paso-green}{PASO}},
  fonttitle=\color{white},
  attach boxed title to top left={yshift=-2mm,xshift=4mm},
  boxed title style={colback=espe-blue,arc=3pt}]
\begin{minipage}[t]{0.48\linewidth}
\textbf{Precondiciones:}
\begin{itemize}[noitemsep,topsep=2pt]
  \item Usuario autenticado.
  \item Transcripci\'on con contenido suficiente.
  \item Servicio IA disponible.
\end{itemize}
\textbf{Pasos de Ejecuci\'on:}
\begin{enumerate}[noitemsep,topsep=2pt]
  \item Abrir una transcripci\'on.
  \item Abrir el chat contextual.
  \item Preguntar algo espec\'ifico sobre la nota.
  \item Esperar respuesta.
  \item Hacer una pregunta de seguimiento.
\end{enumerate}
\end{minipage}\hfill
\begin{minipage}[t]{0.48\linewidth}
\textbf{Resultado Esperado:}
\begin{itemize}[noitemsep,topsep=2pt]
  \item El chat responde sin error HTTP.
  \item La respuesta usa informaci\'on de la nota.
  \item El historial muestra pregunta y respuesta.
\end{itemize}
\vspace{4pt}
\textbf{Resultado Obtenido:}
\begin{itemize}[noitemsep,topsep=2pt]
  \item[$\checkmark$] Responde sin error HTTP.
  \item[$\checkmark$] Respuesta con info de la nota.
  \item[$\checkmark$] Historial visible.
\end{itemize}
\end{minipage}
\end{tcolorbox}

\newpage

% ── 4. RESUMEN ───────────────────────────────────────────────────────────────
\section{Resumen de Ejecuci\'on}

\begin{center}
\begin{tabular}{L{3.5cm} C{2.2cm} C{2.2cm} C{4cm}}
\toprule
\rowcolor{header-bg}
\textcolor{white}{\textbf{Estado}} &
\textcolor{white}{\textbf{Cantidad}} &
\textcolor{white}{\textbf{\% Total}} &
\textcolor{white}{\textbf{Casos}} \\
\midrule
\statusChip{paso-green}{PASO}           & 9 & 90\% & CP-01 al CP-08, CP-10 \\
\rowcolor{row-alt}
\statusChip{pendiente-orange}{PENDIENTE} & 1 & 10\% & CP-09 \\
\statusChip{critica}{FALLO}             & 0 &  0\% & -- \\
\bottomrule
\end{tabular}
\end{center}

\vspace{0.5cm}

\begin{tcolorbox}[colback=row-alt, colframe=espe-blue, boxrule=0.8pt, arc=3pt,
  title={\textbf{Observaciones}}, fonttitle=\color{white},
  attach boxed title to top left={yshift=-2mm,xshift=4mm},
  boxed title style={colback=espe-blue}]
\begin{itemize}
  \item Los flujos cr\'iticos \textbf{CP-03} (Grabaci\'on) y \textbf{CP-04} (Transcripci\'on en vivo) pasaron satisfactoriamente.
  \item \textbf{CP-09} queda pendiente por dependencia de credenciales IA locales; se retomar\'a en el re-testing del D\'ia 3.
  \item Los defectos preliminares BUG-01 a BUG-04 documentados en la secci\'on B.c ser\'an registrados en Jira con evidencia.
\end{itemize}
\end{tcolorbox}

\vspace{0.5cm}

% ── 5. REFERENCIAS ───────────────────────────────────────────────────────────
\section*{Referencias}
\addcontentsline{toc}{section}{Referencias}

\begin{itemize}
  \item \texttt{docs/sqa/sqa-planning.md} -- Plan riguroso SQA Sprint 18--21 Jul 2026.
  \item \texttt{docs/sqa/persona-d-manual-testing.md} -- Fuente de casos y matriz.
  \item IEEE 829-2008 -- Standard for Software and System Test Documentation.
\end{itemize}

\end{document}
